\documentclass[12pt,a4paper]{article}
% Preâmbulo institucional comum — DEE/IFMA
\usepackage[utf8]{inputenc}
\usepackage[T1]{fontenc}
\usepackage[brazil]{babel}
\usepackage{lmodern}
\usepackage{microtype}
\usepackage{geometry}
\usepackage{xcolor}
\usepackage{graphicx}
\usepackage{booktabs}
\usepackage{tabularx}
\usepackage{array}
\usepackage{enumitem}
\usepackage{hyperref}
\usepackage{lastpage}
\usepackage{fancyhdr}
\usepackage{setspace}
\usepackage{titlesec}
\usepackage{tcolorbox}
\usepackage{ragged2e}

\geometry{a4paper,top=2.2cm,bottom=2.2cm,left=2.5cm,right=2.5cm}
\definecolor{azulpetroleo}{HTML}{0F4C5C}
\definecolor{ambar}{HTML}{C8892D}
\definecolor{cinzaclaro}{HTML}{F2F5F7}
\definecolor{cinzaescuro}{HTML}{374151}
\hypersetup{colorlinks=true,linkcolor=azulpetroleo,urlcolor=azulpetroleo,citecolor=azulpetroleo}
\setlength{\parindent}{1.25cm}
\setlength{\parskip}{0.35em}
\onehalfspacing
\titleformat{\section}{\large\bfseries\color{azulpetroleo}}{}{0pt}{}
\titleformat{\subsection}{\normalsize\bfseries\color{azulpetroleo}}{}{0pt}{}
\pagestyle{fancy}
\fancyhf{}
\fancyhead[L]{\small Departamento de Eletroeletrônica — DEE}
\fancyhead[R]{\small IFMA — Campus São Luís/Monte Castelo}
\fancyfoot[C]{\small Página \thepage\ de \pageref{LastPage}}
\renewcommand{\headrulewidth}{0.4pt}
\renewcommand{\footrulewidth}{0.4pt}
\newcommand{\campo}[1]{\textcolor{ambar}{\textbf{[#1]}}}
\newcommand{\instituicao}{Instituto Federal de Educação, Ciência e Tecnologia do Maranhão}
\newcommand{\campus}{Campus São Luís/Monte Castelo}
\newcommand{\departamento}{Departamento de Eletroeletrônica — DEE}
\newcommand{\cidade}{São Luís — MA}
\newcommand{\cabecalhoInstitucional}{%
\begin{center}
{\Large\bfseries\color{azulpetroleo}\instituicao}\\[2mm]
{\large\bfseries\campus}\\[1mm]
{\normalsize\departamento}
\end{center}
\vspace{2mm}
\hrule
\vspace{4mm}
}
\newtcolorbox{destaque}{colback=cinzaclaro,colframe=azulpetroleo,boxrule=0.6pt,arc=2mm}

\begin{document}
\cabecalhoInstitucional
\begin{center}
{\LARGE\bfseries\color{azulpetroleo} Programa de Modernização Tecnológica dos Laboratórios do DEE}\\[2mm]
{\large Prospecção de Doações, Parcerias e Cooperação Tecnológica}
\end{center}

\section{Apresentação}
O Departamento de Eletroeletrônica do IFMA — Campus São Luís/Monte Castelo apresenta o presente Programa de Modernização Tecnológica com o objetivo de ampliar e atualizar a infraestrutura destinada ao ensino, à pesquisa, à extensão e à formação profissional nas áreas de Eletrotécnica, Eletrônica, Automação e Engenharia Elétrica.

\section{Objetivos}
\begin{itemize}[leftmargin=*]
\item modernizar laboratórios e bancadas didáticas;
\item aproximar a formação acadêmica das tecnologias utilizadas no setor produtivo;
\item ampliar a capacidade de ensino prático, pesquisa aplicada e extensão;
\item estabelecer parcerias permanentes com empresas do setor elétrico, eletrônico, de automação, telecomunicações e tecnologia;
\item estimular projetos de inovação, eficiência energética, Indústria 4.0, IoT, proteção, medição e qualidade da energia.
\end{itemize}

\section{Público beneficiado}
Cursos atendidos: Técnico em Eletrotécnica, Técnico em Eletrônica, Técnico em Automação e Engenharia Elétrica, além de projetos de pesquisa, extensão, iniciação científica, TCCs e capacitações institucionais.

\section{Áreas prioritárias}
\begin{enumerate}[leftmargin=*]
\item Automação Industrial e Controle;
\item Máquinas Elétricas e Acionamentos;
\item Instalações, Medidas e Qualidade da Energia;
\item Sistemas de Energia e Proteção;
\item Eletrônica, Instrumentação e IoT;
\item Redes Industriais e Telecomunicações;
\item Computação aplicada à Engenharia.
\end{enumerate}

\section{Modalidades de apoio desejadas}
\begin{itemize}[leftmargin=*]
\item doação de equipamentos novos, de demonstração ou descontinuados, desde que plenamente funcionais;
\item cessão de kits didáticos, licenças e softwares educacionais;
\item estruturação ou modernização de laboratórios;
\item treinamento de docentes e discentes;
\item apoio financeiro para aquisição de equipamentos;
\item cooperação técnica em projetos de pesquisa, extensão e inovação.
\end{itemize}

\section{Empresas prioritárias}
A prospecção inicial contempla, entre outras, Schneider Electric, Vale, Eneva, Equatorial Energia, WEG, Rockwell Automation, Siemens, Altus, Intelbras, Petrobras, Festo, SMC, ABB, Alcoa/Alumar, Suzano, Hitachi Energy, Eaton, Phoenix Contact, Emerson, Fluke, Keysight, Tektronix/Keithley, National Instruments, Yaskawa, Mitsubishi Electric, Bosch Rexroth, Cisco, Huawei, Dell Technologies e Lenovo.

\section{Contrapartidas institucionais possíveis}
\begin{itemize}[leftmargin=*]
\item divulgação institucional da parceria, observadas as normas do IFMA;
\item identificação do laboratório apoiado;
\item utilização dos equipamentos em atividades didáticas e de pesquisa;
\item relatórios de impacto e uso dos equipamentos;
\item ações de capacitação e demonstração tecnológica;
\item participação da empresa em eventos acadêmicos e tecnológicos do campus.
\end{itemize}

\section{Governança e formalização}
Toda doação ou parceria será formalizada por processo administrativo, com análise dos setores competentes do IFMA e observância da legislação e das normas vigentes aplicáveis à Administração Pública Federal.

\begin{destaque}
\textbf{Coordenação institucional:} Departamento de Eletroeletrônica — DEE\\
\textbf{Responsável:} \campo{Nome do responsável}\\
\textbf{Função:} \campo{Chefe do DEE / responsável pelo projeto}\\
\textbf{Contato:} \campo{e-mail institucional / telefone}
\end{destaque}

\vfill
\begin{center}
\cidade, \campo{data}.\\[15mm]
\rule{8cm}{0.4pt}\\
\campo{Nome do responsável}\\
\campo{Cargo/Função}
\end{center}
\end{document}
